% ============================================================
% f2014_sichuan_science_q05_flowchart.tex
% 2014 四川卷（理科）第 5 题：程序框图
% 来源 PDF：普通高考/2014/2014四川理.pdf
% 原页：第 1 页（A3 横排）左栏下部，印面 919，第 5 题题干右下
% 结构：开始 → 输入 x,y → 判断「x≥0，y≥0，x+y≤1？」（菱形）；
%       是：S=2x+y；否：S=1；两路汇合 → 输出 S → 结束。
% ============================================================
\begin{tikzpicture}[>=Stealth, line width=0.7pt, font=\small,
  x=1cm, y=1cm, line cap=round, line join=round]

  % 开始
  \draw[rounded corners=6pt] (-0.95,0.45) rectangle (0.95,-0.45);
  \node at (0,0) {开始};
  \draw[->] (0,-0.45) -- (0,-0.95);

  % 输入
  \draw (-1.25,-0.95) -- (1.75,-0.95) -- (1.25,-1.65) -- (-1.75,-1.65) -- cycle;
  \node at (0,-1.30) {输入 $x$，$y$};
  \draw[->] (0,-1.65) -- (0,-2.15);

  % 判断（菱形，与原页一致）
  \draw (-2.95,-3.15) -- (0,-2.15) -- (2.95,-3.15) -- (0,-4.15) -- cycle;
  \node at (0,-2.95) {$x\geqslant 0$，$y\geqslant 0$，};
  \node at (0,-3.38) {$x+y\leqslant 1$？};

  % 是：左支
  \draw[->] (-2.95,-3.15) -- (-4.15,-3.15) -- (-4.15,-4.30);
  \node at (-3.35,-2.83) {是};
  \draw (-5.45,-4.30) rectangle (-2.85,-5.10);
  \node at (-4.15,-4.70) {$S=2x+y$};

  % 否：右支
  \draw[->] (2.95,-3.15) -- (4.15,-3.15) -- (4.15,-4.30);
  \node at (3.35,-2.83) {否};
  \draw (2.85,-4.30) rectangle (5.45,-5.10);
  \node at (4.15,-4.70) {$S=1$};

  % 两路汇合
  \draw (-4.15,-5.10) -- (-4.15,-5.50) -- (4.15,-5.50) -- (4.15,-5.10);
  \draw[->] (0,-5.50) -- (0,-5.85);

  % 输出
  \draw (-1.15,-5.85) -- (1.85,-5.85) -- (1.35,-6.55) -- (-1.65,-6.55) -- cycle;
  \node at (0.10,-6.20) {输出 $S$};
  \draw[->] (0.10,-6.55) -- (0.10,-6.95);

  % 结束
  \draw[rounded corners=6pt] (-0.85,-6.95) rectangle (1.05,-7.85);
  \node at (0.10,-7.40) {结束};
\end{tikzpicture}
